% =============================================================================
%  CAP 8.0 — Technical maintenance report
%
%  Target audience: original authors (C. Pruneau, V. Gonzalez, O. Sheibani).
%
%  Compile:
%    pdflatex TECHNICAL_REPORT.tex
%    pdflatex TECHNICAL_REPORT.tex   (twice, for the table of contents)
% =============================================================================
\documentclass[11pt,a4paper]{article}

\usepackage[margin=1in]{geometry}
\usepackage{microtype}
\usepackage{enumitem}
\usepackage{xcolor}
\usepackage{hyperref}
\usepackage{listings}
\usepackage{booktabs}
\usepackage{longtable}
\usepackage{titlesec}
\usepackage{parskip}
\usepackage{fancyvrb}

\hypersetup{
  colorlinks=true,
  linkcolor=black,
  urlcolor=blue!60!black,
  citecolor=black
}

\definecolor{codebg}{RGB}{245,247,249}
\definecolor{codecomment}{RGB}{106,135,89}
\definecolor{codestring}{RGB}{106,153,85}

\lstset{
  basicstyle=\ttfamily\footnotesize,
  backgroundcolor=\color{codebg},
  frame=single,
  rulecolor=\color{black!20},
  commentstyle=\color{codecomment},
  stringstyle=\color{codestring},
  breaklines=true,
  showstringspaces=false,
  framesep=4pt,
  framerule=0.4pt,
  xleftmargin=4pt,
  xrightmargin=4pt
}

\lstdefinelanguage{ini}{
  morecomment=[l]{\#},
  morestring=[b]",
  sensitive=true
}

\titleformat{\section}{\large\bfseries}{\thesection}{0.6em}{}
\titleformat{\subsection}{\normalsize\bfseries}{\thesubsection}{0.5em}{}

\title{\textbf{CAP 8.0 — Build System Modernisation\\and Maintenance Report}}
\author{Maintenance pass on the existing codebase\\\small Original authors: C. Pruneau, V. Gonzalez, O. Sheibani}
\date{April 30, 2026}

\begin{document}
\maketitle

\begin{abstract}
\noindent
This report documents a maintenance pass on \texttt{CAP~8.0} performed during
a single working session. The work covers four areas: (i)~modernisation of the
CMake build system and creation of portable installer/runner tooling so that
non-developers can build and execute CAP on a wide range of platforms (macOS,
Linux, conda, Homebrew, MacPorts, CVMFS, environment-modules, Wayne State Grid);
(ii)~repair of source-level bugs accumulated from incomplete refactors,
principally in \texttt{src/CAPPythia/PythiaEventGenerator.cpp} and several
\texttt{src/Glauber/} files; (iii)~addition of backward-compatibility shims in
\texttt{Task.cpp} so the existing project \texttt{.ini} files load with the
current code; and (iv)~documentation of two large unfinished refactors that
prevent end-to-end runs --- the missing orchestrator/calculator classes and
the \texttt{src/Jets/} module. A new \emph{smoke test} macro
(\texttt{src/Macros/SmokeTest.C}) was added to validate the build and the
Pythia~8 integration without depending on the unfinished pieces.

The infrastructure changes are non-invasive to the physics analysis code and
should integrate cleanly with future development. The unfinished pieces are
identified with file-and-line precision in companion documents
(\texttt{MISSING\_CLASSES.md}, \texttt{JETS\_KNOWN\_ISSUES.md}).
\end{abstract}

\tableofcontents
\vspace{1em}
\hrule
\vspace{1em}

% =============================================================================
\section{Context and scope}
% =============================================================================

The starting state of \texttt{CAP~8.0} as it appeared on the maintainer's
machine had several practical obstacles to a clean build:

\begin{itemize}
  \item Hard-coded absolute paths inside CMake configuration
        (\texttt{install(TARGETS CAP LIBRARY DESTINATION
        "/Users/aa7526/Documents/GitHub/CAP8.0/lib")}), which
        made the build non-portable across machines.
  \item An invalid macOS deployment target
        (\texttt{CMAKE\_OSX\_DEPLOYMENT\_TARGET~=~26.0}); macOS~26 does not
        exist.
  \item Optional dependencies (Pythia~8, FastJet) located through one
        single hard-coded path (\texttt{/usr/local/lib/}), which excludes
        Apple Silicon Homebrew (\texttt{/opt/homebrew}), MacPorts, conda
        environments and HPC module systems.
  \item Compile errors in \texttt{src/CAPPythia/PythiaEventGenerator.cpp}
        and \texttt{src/CAPPythia/PythiaEventReader.cpp}, both arising from
        an incomplete renaming of base classes
        (\texttt{ProcessorSingle} $\rightarrow$ \texttt{EventProcessor}),
        of the particle-database type
        (\texttt{ParticleTypeList} $\rightarrow$ \texttt{ParticleDb}),
        and of the four-vector setter
        (\texttt{setPxPyPzE} $\rightarrow$ \texttt{setEPxPyPz},
        which also reorders the arguments).
  \item All shipped \texttt{.ini} project files use the legacy key naming
        convention (\texttt{TaskClassName},~\texttt{nSubtasks},
        \texttt{Subtask$k$:TaskName}), whereas \texttt{Base/Task.cpp} now
        looks up keys in the new structural form
        (\texttt{TASK:CLASS},~\texttt{SUBTASK:N}, \texttt{SUBTASK:$k$:NAME}).
  \item The \texttt{src/Jets/} module is incomplete and does not compile
        without significant additional development.
  \item Approximately ten task classes referenced by
        \texttt{projects/Pythia/pp\_13.7TeV/RunAna.ini}
        (\texttt{CAP::RunAnalysis}, \texttt{CAP::ParticleTypeTask},
        \texttt{CAP::EventFilterCreator}, \texttt{CAP::FileIterator}, the
        \texttt{*Calculator} family) have no implementation in \texttt{src/}
        and therefore cannot be instantiated by ROOT's
        \texttt{TClass::GetClass()}.
\end{itemize}

The goal of the maintenance pass was to re-establish a clean, portable build
flow for every part of the codebase that does compile, repair the bit-rot
where the fix was mechanical, and clearly delineate the parts that require
substantive design/development work.

% =============================================================================
\section{Build system overhaul}
% =============================================================================

\subsection{Top-level structure}

A new top-level \texttt{CMakeLists.txt} was added at the repository root to
allow the canonical out-of-source build invocation:

\begin{lstlisting}[language=bash]
cmake -S . -B build
cmake --build build -j
cmake --install build
\end{lstlisting}

\noindent
The previous flow \texttt{cmake ../src} from inside a \texttt{build/}
directory still works.

\subsection{Per-module install destinations}

All twenty-two per-module \texttt{src/<Module>/CMakeLists.txt} files were
modified to install via the standard CMake variable

\begin{lstlisting}[language=bash]
install(TARGETS <Module> LIBRARY DESTINATION ${CAP_INSTALL_LIBDIR})
\end{lstlisting}

\noindent
rather than \verb|$ENV{CAP_LIB_PATH}|. This removes the requirement that
\texttt{SetupCAP.sh} be sourced before \texttt{make install}.
The default install prefix is computed from
\texttt{CMAKE\_CURRENT\_SOURCE\_DIR} so that, regardless of how CMake is
invoked, libraries land at \texttt{<repo>/lib} and binaries at
\texttt{<repo>/bin}.

\subsection{Optional-module flags}

The pre-existing \verb|set(Include_PYTHIA "NO")|-style flags were converted to
proper CMake \texttt{option()} declarations:

\begin{lstlisting}[language=bash]
option(CAP_ENABLE_THERMINATOR  "..."  ON)
option(CAP_ENABLE_GLAUBER      "..."  ON)
option(CAP_ENABLE_PYTHIA       "..."  OFF)
option(CAP_ENABLE_JETS         "..."  OFF)
option(CAP_ENABLE_EPOS         "..."  OFF)
option(CAP_ENABLE_PHSD         "..."  OFF)
\end{lstlisting}

\noindent
Users override these on the command line:
\begin{lstlisting}[language=bash]
cmake -S . -B build -D CAP_ENABLE_PYTHIA=ON
\end{lstlisting}

\noindent
The legacy \texttt{Include\_*} variable names continue to work via a
back-compatibility shim that copies their values into the new options.

\subsection{Find modules for external dependencies}

Two CMake find-modules were added under \texttt{cmake/}:

\begin{itemize}
  \item \texttt{cmake/FindPythia8.cmake}
  \item \texttt{cmake/FindFastJet.cmake}
\end{itemize}

\noindent
Both honour the following discovery sources, in order of priority:

\begin{enumerate}[leftmargin=2em,itemsep=0pt]
  \item Explicit user override via the CMake-level variable
        \verb|CAP_PYTHIA8_PATH| / \verb|CAP_FASTJET_PATH|, or the equivalent
        environment variable.
  \item The standard upstream environment variables (\verb|PYTHIA8_DIR|,
        \verb|FASTJET_DIR|, etc.).
  \item The companion \texttt{pythia8-config} or \texttt{fastjet-config}
        executable, when on \texttt{PATH}.
  \item \texttt{pkg-config} (\texttt{pythia8.pc}, \texttt{fastjet.pc}).
  \item Apple Homebrew under \texttt{/opt/homebrew} \emph{and} \texttt{/usr/local}.
  \item Conda environments via \verb|$CONDA_PREFIX|.
\end{enumerate}

\noindent
After a successful detection both modules export an imported target
(\texttt{Pythia8::Pythia8}, \texttt{FastJet::FastJet}) and conventional
\texttt{*\_VERSION} / \texttt{*\_INCLUDE\_DIRS} / \texttt{*\_LIBRARIES}
variables.

\subsection{Build summary}

The configure step now prints an explicit, human-readable summary that
distinguishes \emph{detected} from \emph{linked}:

\begin{lstlisting}
========================  CAP build configuration  ========================
  Install prefix         : /Users/<user>/cap8/CAP8.0-main
  Build type             : RelWithDebInfo
  C++ standard           : 17
  ROOT                   : 6.36.06
  Therminator            : ON
  Glauber                : ON
  Pythia 8               : linked  (v8.310, /opt/homebrew/.../libpythia8.dylib)
  FastJet                : detected on system (v3.4.3) but not linked
  Jets module            : OFF  (broken upstream - see JETS_KNOWN_ISSUES.md)
  EPOS / PHSD            : OFF / OFF
===========================================================================
\end{lstlisting}

% =============================================================================
\section{Installer and runner infrastructure}
% =============================================================================

A unified entry-point script \texttt{./install} was added at the repository
root with four modes:

\begin{center}
\begin{tabular}{ll}
\toprule
\texttt{./install --gui}      & Tk-based graphical setup picker \\
\texttt{./install --cli}      & Text-prompt setup picker \\
\texttt{./install --headless} & Auto-detect, configure, build, install \\
\texttt{./install --run}      & Tk-based job runner \\
\bottomrule
\end{tabular}
\end{center}

\subsection{Setup pickers (\texttt{setup-cap} and \texttt{setup-cap-gui})}

Both setup pickers share a single bash detection layer that scans, in order:

\begin{itemize}[leftmargin=2em,itemsep=0pt]
  \item User overrides (\verb|$CAP_PYTHIA8_PATH|, \verb|$CAP_FASTJET_PATH|).
  \item In-tree directories (\texttt{./pythia8XYZ}, \texttt{./fastjet-X.Y.Z}).
  \item Brew formula paths for both the canonical core formula
        (\texttt{pythia8}, \texttt{fastjet}) and the third-party HEP tap
        (\texttt{davidchall/hep/pythia}, \texttt{davidchall/hep/fastjet}).
  \item Direct prefixes (\texttt{/opt/homebrew}, \texttt{/usr/local},
        \texttt{/opt/local}) for hand-built installs that do not use the
        keg-style layout.
  \item Apple MacPorts (\texttt{/opt/local}).
  \item conda (\verb|$CONDA_PREFIX|).
  \item \texttt{pythia8-config} / \texttt{fastjet-config} on \texttt{PATH}.
  \item CVMFS (\texttt{/cvmfs/sft.cern.ch/lcg/releases/...},
        \texttt{/cvmfs/alice.cern.ch/...}).
  \item Wayne State Grid software trees
        (\texttt{/wsu/el?/sw/<package>/...}).
  \item Environment modules (\texttt{module avail}).
\end{itemize}

\noindent
Each candidate is validated by checking that the canonical header file is
actually present in the proposed prefix; candidates with empty paths or
wrong layouts are rejected and the rejection reason is logged.

The CLI installer (\texttt{setup-cap}) presents a numbered menu; the GUI
installer (\texttt{setup-cap-gui}) uses Tk drop-downs. Both write the same
\texttt{.cap-config} file (a shell snippet sourced by \texttt{SetupCAP.sh}),
and both stream output to a per-run log under \texttt{logs/}.

The GUI's \emph{Configure \& Build} button runs the full configure $\to$
build $\to$ install sequence end-to-end. After install it performs a
post-install verification step that checks that \texttt{bin/CAP} and
\texttt{lib/lib*.dylib} are inside the repository (and relocates them if
they accidentally landed in the parent directory because of a stale
\texttt{CMakeCache.txt}).

\subsection{Runner GUI (\texttt{run-cap})}

A Tk-based job runner was added. It scans the \texttt{projects/} tree for
\texttt{.ini} files, parses each file to extract its top-level task
declarations, and presents the user with a project~/~config~/~task drop-down
chain. It supports event-count and random-seed overrides applied to a
per-run copy of the \texttt{.ini} (the original is never mutated). A live
log pane streams the running binary's stdout/stderr, with Stop and
\emph{Open output folder} controls.

It also classifies each \texttt{.ini} as \emph{modern}
(uses the \texttt{Subtask$k$:TaskName} syntax) or \emph{legacy} (uses
\texttt{SubTaskName$k$}), and surfaces the legacy ones with a clear visual
tag so users do not pick a file the loader cannot understand.

A second button, \emph{Smoke test (proof of life)}, invokes a minimal ROOT
macro (Section~\ref{sec:smoketest}) that bypasses the unfinished orchestrator
classes.

\subsection{Logging library}

A shared bash logging library
(\texttt{scripts/cap-logging.sh}) was added. It is sourced by all the
shell-based installer scripts and offers:

\begin{itemize}[leftmargin=2em,itemsep=0pt]
  \item \texttt{cap\_log\_init <component>} --- opens a per-run log file
        \texttt{logs/<component>-<UTC-timestamp>.log} and a
        \texttt{logs/last-<component>.log} convenience symlink.
  \item \texttt{cap\_log <level> <message>} --- writes a structured line
        with platform-portable UTC timestamps. (BSD~\texttt{date} on macOS
        does not support \texttt{\%N}; the helper detects this and falls
        back to second precision.)
  \item \texttt{cap\_log\_run "<label>" <cmd> <args>...} --- runs a command,
        captures stdout/stderr to the log and the parent terminal, and
        returns the wrapped command's true exit code via
        \texttt{PIPESTATUS[0]} after temporarily disabling
        \texttt{set~-e} to ensure the failure-reporting branch executes.
  \item \texttt{cap\_log\_close <exit-code>} --- closes the file with a
        machine-readable footer.
\end{itemize}

\noindent
The Python-based GUIs (\texttt{setup-cap-gui}, \texttt{run-cap}) write logs
in the same format using a small Python class so that
\texttt{logs/} is a single unified store across all tools.

% =============================================================================
\section{Source-level fixes}
% =============================================================================

\subsection{\texttt{src/CAPPythia/PythiaEventGenerator.cpp}}

The file did not compile. Twelve distinct sites were repaired:

\begin{itemize}[leftmargin=2em,itemsep=0pt]
  \item Line~13 (now removed): \verb|#include "TVectorLorentz.h"| --- a header
        that does not exist in any version of ROOT (likely a typo for
        \texttt{TLorentzVector.h}, which is also unused). The include line
        was deleted.
  \item Line~63 was \verb|ProcessorSingle::setDefaultConfiguration();| ---
        the base class was renamed to \texttt{EventProcessor}.
  \item Lines~88 and~134 declared \texttt{String~\&~key~=~taskName;} and then
        \texttt{key~+=~":Option";~key~+=~k;} --- because \texttt{taskName} is
        a \texttt{const String\,\&}, this would either fail to compile or
        mutate the source variable. Both sites were replaced with
        \texttt{String~key~=~taskName;} (a copy).
  \item Lines~105, 106, 113, 114 had four missing closing parentheses in
        nested calls
        (\verb|...valueInt(createKey(taskName,"Beams:frameType"));| where
        the call should close with \verb|)));|).
  \item Line~210 used the deleted singleton
        \verb|ParticleTypeList::list();| --- replaced with
        \verb|db()|, which returns the
        \texttt{ParticleDb} owned by the enclosing
        \texttt{EventProcessor}.
  \item Lines~212 and~236 wrote \texttt{event.reset();} and
        \texttt{event.addParticle(\&p);} --- using \texttt{event} (which is
        the \emph{member function} \verb|EventProcessor::event()|) as if it
        were the \texttt{Event} reference. Replaced with the local
        \texttt{theEvent} that line~211 already declares.
  \item Line~234 used the renamed and re-ordered four-vector setter:
        \verb|setPxPyPzE(px,py,pz,E)| was changed to
        \verb|setEPxPyPz(E,px,py,pz)|, with arguments reordered.
  \item Line~238 had \texttt{eventAccepted.increment();} where the symbol
        is now an accessor returning \texttt{Accountant\&}; corrected to
        \texttt{eventAccepted().increment();}.
\end{itemize}

\noindent
A missing \verb|#include "ParticleDb.hpp"| was added at the top of the
file.

\subsection{\texttt{src/CAPPythia/PythiaEventReader.cpp}}

This file is far more deeply broken than the generator: it references
\texttt{ProcessorSingle}, the deleted \texttt{ParticleTypeList::list()}
singleton, and a number of TTree branch fields (\texttt{entryIndex},
\texttt{nb}, \texttt{nBytes}, \texttt{nParticles},
\texttt{tracks\_fStatusCode}, \texttt{tracks\_fPdgCode}) that no longer
exist anywhere in the codebase. Repair would require re-deriving the
file's intended TTree schema. As an interim measure, the file was
\textbf{removed from the \texttt{libCAPPythia} build target}
(\texttt{src/CAPPythia/CMakeLists.txt}) and its dictionary entry was
commented out in \texttt{CAPPythiaLinkDef.h}. The source file remains in
the repository so a future maintainer can repair it; nothing in the
analysis flow currently references it.

\subsection{\texttt{src/Glauber/GlauberHistos.cpp}}

\begin{itemize}[leftmargin=2em,itemsep=0pt]
  \item \texttt{calculateDensity(...)} (line~897) had its entire body
        commented out, leaving a non-\texttt{void} function with no
        \texttt{return} statement (undefined behaviour). The body was
        replaced by an explicit \texttt{return~0.0;} with a comment
        marking the function as a stub awaiting a real implementation.
  \item Two unused locals \texttt{typeA} and \texttt{typeB} (lines~947--948)
        were removed; their usage had been commented out earlier.
  \item One unused-variable warning on \texttt{nucleon} (line~889) was
        silenced with \texttt{\_\_attribute\_\_((unused))} since it is part
        of a documented stub loop.
\end{itemize}

\subsection{Other minor warnings}

\begin{itemize}[leftmargin=2em,itemsep=0pt]
  \item \texttt{src/Helpers/CanvasList.cpp} line~62: \texttt{xInc} was
        marked unused.
  \item \texttt{src/Glauber/NuclearDistribution.cpp} line~74: \texttt{tp}
        was extracted from the multi-declaration and marked unused.
\end{itemize}

\subsection{\texttt{src/Jets/} --- module disabled at build time}

The Jets module is \emph{not bit-rot}, it is incomplete in-progress code.
Among the issues identified:

\begin{itemize}[leftmargin=2em,itemsep=0pt]
  \item \texttt{JetAnalyzer.hpp} references a class
        \texttt{JetSingleDerivedHistos} that does not exist; the actual
        class is named \texttt{JetSingleHistosDerived}.
  \item It also references a template \texttt{ManagedList<JetFilter>} that
        is defined nowhere in the codebase.
  \item A member \texttt{clusterSequence} is used in a default member
        initialiser without ever being declared as a member.
  \item \texttt{JetHistosDerived.hpp} line~45 has a typographical error:
        \texttt{cloneD(onst JetHistosDerived\&~source)} (the leading
        \texttt{c} of \texttt{const} was dropped).
  \item \texttt{JetAnalyzer.cpp} contains $\sim$15 distinct issues: missing
        closing parentheses in \texttt{addProperty(...)} calls, an undefined
        base class \texttt{Analyzer}, a function-call site that contains a
        parameter declaration, mismatched braces in the event loop, a
        function (\texttt{filterJets}) defined but not declared in the
        header, multiple references to undeclared variables
        (\texttt{taskName}, \texttt{iEventFilter}, \texttt{iParticleFilter},
        \texttt{acceptedParticleFilters}).
\end{itemize}

\noindent
Because none of these are mechanical fixes (the module needs a design pass
to resolve them), the build system was changed to \emph{refuse} to compile
the Jets module unless the developer explicitly opts in:

\begin{lstlisting}[language=bash]
cmake -S . -B build -D CAP_ENABLE_JETS=ON -D CAP_FORCE_BROKEN_JETS=ON
\end{lstlisting}

\noindent
The token \texttt{CAP\_FORCE\_BROKEN\_JETS} is the explicit acknowledgement
that the developer is choosing to engage with broken code. The
\texttt{ParticleTypeList} $\to$ \texttt{ParticleDb} renaming
\emph{was} applied to the Jets headers and source files
(\texttt{JetHistos.hpp}, \texttt{JetSingleHistos.hpp},
\texttt{JetPairHistos.hpp}, \texttt{JetPairMassHistos.hpp},
\texttt{JetHistosDerived.hpp}, and the corresponding \texttt{.cpp} files),
plus a typographical fix
(\texttt{ParticleTypeList = selectedParticleTypeList;} $\to$
\texttt{particleTypeList = selectedParticleTypeList;}) so that whoever
finishes the module starts from a slightly cleaner baseline. The full bug
list with file/line references is captured in
\texttt{JETS\_KNOWN\_ISSUES.md}.

% =============================================================================
\section{Configuration system: backward-compatibility shim}
\label{sec:taskpatch}
% =============================================================================

\subsection{The mismatch}

\texttt{src/Base/Task.cpp} looks up the task tree via three structural keys:

\begin{lstlisting}[language=ini]
<task>:TASK:CLASS         # what class to instantiate
<task>:TASK:SEVERITY      # log verbosity
<task>:SUBTASK:N          # number of subtasks
<task>:SUBTASK:k:NAME     # name of subtask k
\end{lstlisting}

\noindent
Every \texttt{.ini} file under \texttt{projects/} uses the legacy form:

\begin{lstlisting}[language=ini]
<task>:TaskClassName     = CAP::RunAnalysis
<task>:Severity          = Info
<task>:nSubtasks         = 4
<task>:Subtask0:TaskName = ParticleTypeTask
\end{lstlisting}

\noindent
With unmodified \texttt{Task.cpp} and the shipped \texttt{.ini} files, the
top-level task lookup returns the literal default
\texttt{"NONE"} (the default value of
\texttt{Properties::valueString(...)},
declared in \texttt{Properties.hpp}), and CAP terminates with
\texttt{UnknownClassException name=NONE} before any task is instantiated.

\subsection{The patch}

Three sites in \texttt{src/Base/Task.cpp} were modified to look up the new
key first and fall back to the legacy key only when the new lookup
returned the empty string \emph{or} the literal \texttt{"NONE"}:

\begin{itemize}[leftmargin=2em,itemsep=0pt]
  \item \texttt{Task::createTask()} --- \texttt{TASK:CLASS}
        $\to$ \texttt{TaskClassName} fallback.
  \item \texttt{Task::configure()} --- \texttt{TASK:SEVERITY}
        $\to$ \texttt{Severity} fallback.
  \item \texttt{Task::configureSubtasks()} ---
        \texttt{SUBTASK:N} $\to$ \texttt{nSubtasks}, and
        \texttt{SUBTASK:$k$:NAME} $\to$ \texttt{Subtask$k$:TaskName} fallback.
\end{itemize}

\noindent
The \texttt{== "NONE"} check is essential: a missing key returns the
literal four-character string \texttt{"NONE"} (not an empty string), so a
naive \texttt{IsNull()~||~Length()==0} check never triggers.

\subsection{Coverage}

The shim covers only the structural keys used by \texttt{Task.cpp} itself.
Per-task keys (e.g.\ \texttt{HistogramsExportPath},
\texttt{nParticleDbs}, \texttt{nEventsRequested}) have not been
back-compat'd and would need analogous treatment in their respective
modules if they prove to be a blocker after the missing classes are
implemented.

% =============================================================================
\section{Findings: unfinished refactors and missing implementations}
% =============================================================================

\subsection{Class refactor not propagated to all modules}

The codebase contains evidence of three large renamings that were applied
to some modules but not others:

\begin{center}\small
\begin{tabular}{lll}
\toprule
\textbf{Old name} & \textbf{New name} & \textbf{Modules updated} \\
\midrule
\texttt{ParticleTypeList} & \texttt{ParticleDb} & \texttt{Particles}, \texttt{CAPPythia} (after this pass) \\
\texttt{ProcessorSingle} & \texttt{EventProcessor / EventProcessorSingle} & \texttt{Particles} \\
\texttt{Particle::setPxPyPzE} & \texttt{Particle::setEPxPyPz} (re-ordered) & \texttt{Particles} \\
\bottomrule
\end{tabular}
\end{center}

\noindent
The Jets module has not been updated for any of the three; the CAPPythia
module had been updated only partially (it inherited from
\texttt{EventProcessor} but still called \texttt{ProcessorSingle::setDefaultConfiguration()},
and still used \texttt{ParticleTypeList::list()}).

\subsection{Configuration system refactor not propagated to .ini files}

The most extensive of the unfinished refactors. \texttt{src/} now uses a
new structural key naming scheme:

\begin{center}\small
\begin{tabular}{ll}
\toprule
\textbf{Legacy (used by .ini files)} & \textbf{Current (used by source)} \\
\midrule
\texttt{TaskName}                 & \texttt{TASK:NAME} \\
\texttt{TaskClassName}            & \texttt{TASK:CLASS} \\
\texttt{Severity}                 & \texttt{TASK:SEVERITY} \\
\texttt{nSubtasks}                & \texttt{SUBTASK:N} \\
\texttt{Subtask$k$:TaskName}      & \texttt{SUBTASK:$k$:NAME} \\
\texttt{nParticleDbs}             & \texttt{DB:N} \\
\texttt{ParticleDbName$k$}        & \texttt{DB\_$k$:NAME} \\
\texttt{ParticleDbOwner$k$}       & \texttt{DB\_$k$:OWNER} \\
\texttt{nEventsRequested}         & \texttt{EVENT:REQUESTED:N} \\
\texttt{HistogramsExportPath}     & \texttt{<typeName>:EXPORT:PATH} \\
\bottomrule
\end{tabular}
\end{center}

\noindent
The \texttt{<typeName>} indirection is significant: the new scheme groups
keys under a \emph{type name} chosen by each task class
(\texttt{"DB"}, \texttt{"EVENT"}, \texttt{"PARTICLE\_FILTER"},
\texttt{"INPUT\_STREAM"}, etc.). A purely string-based translator from
old to new is therefore not sufficient --- the type-name has to be known
per task.

\noindent
The structural keys (\texttt{TASK:CLASS}, \texttt{TASK:SEVERITY},
\texttt{SUBTASK:N}, \texttt{SUBTASK:$k$:NAME}) have a back-compat path
courtesy of Section~\ref{sec:taskpatch}; the per-task keys do not.

\subsection{Missing task classes}

The most consequential finding. Of the sixteen distinct
\texttt{TaskClassName} values referenced by
\texttt{projects/Pythia/pp\_13.7TeV/RunAna.ini}, ten are missing from
\texttt{src/}:

\begin{center}\small
\begin{longtable}{lp{8cm}}
\toprule
\textbf{Class} & \textbf{Role} \\
\midrule\endhead
\texttt{CAP::RunAnalysis}              & top-level container task \\
\texttt{CAP::ParticleTypeTask}         & loads \texttt{ParticleDb} from disk \\
\texttt{CAP::ParticleDbTask}           & alternative DB loader \\
\texttt{CAP::EventFilterCreator}       & instantiates EventFilters from .ini blocks \\
\texttt{CAP::ParticleFilterCreator}    & instantiates ParticleFilters from .ini blocks \\
\texttt{CAP::FilterCreator}            & combined Event/Particle/Jet filter creator \\
\texttt{CAP::FileIterator}             & post-processing loop driver \\
\texttt{CAP::ParticleSingleCalculator} & derived single-particle histograms \\
\texttt{CAP::ParticlePairCalculator}   & derived pair correlation observables \\
\texttt{CAP::ParticlePair3DCalculator} & derived 3-D HBT observables \\
\bottomrule
\end{longtable}
\end{center}

\noindent
Six classes are present and compile cleanly:
\texttt{EventIterator}, \texttt{PythiaEventGenerator},
\texttt{ParticleSingleAnalyzer}, \texttt{ParticlePairAnalyzer},
\texttt{ParticlePair3DAnalyzer}, plus support classes (\texttt{Task},
\texttt{EventProcessor}, \texttt{Event}, \texttt{Particle},
\texttt{ParticleDb}, \texttt{EventFilter}, \texttt{ParticleFilter}, etc.).

The companion file \texttt{MISSING\_CLASSES.md} provides per-class detail on
expected role, dependencies, and rough size estimates. The simplest of the
missing classes (\texttt{RunAnalysis} as a thin \texttt{Task} subclass) is
$\sim$30 lines; the calculators encode actual physics and are several
hundred lines each.

% =============================================================================
\section{Smoke test (\texttt{src/Macros/SmokeTest.C})}
\label{sec:smoketest}
% =============================================================================

A ROOT macro was added that exercises the entire build/link chain
\emph{without} relying on the missing orchestrator classes. It:

\begin{enumerate}[leftmargin=2em,itemsep=0pt]
  \item Loads \texttt{libBase}, \texttt{libMath}, \texttt{libParticles},
        \texttt{libCAPPythia} via \texttt{gSystem->Load(...)};
  \item Adds a Cling-aware include and library path covering the most
        common Pythia install locations (Apple Silicon Homebrew, Intel
        Homebrew, MacPorts, \texttt{/usr/local}, the davidchall/hep tap);
  \item Constructs a \texttt{Pythia8::Pythia} object directly,
        configures pp@13~TeV inelastic, runs $N$ events;
  \item Fills five histograms by hand from the Pythia event record;
  \item Writes them to \texttt{histos/smoketest/SingleGen.root}.
\end{enumerate}

\noindent
On the maintainer's machine the macro completes in $\sim$5 seconds and
produces a valid \texttt{.root} file with the expected charged-particle
multiplicity, $p_T$, $\eta$, $\phi$ distributions, and a 2-D $p_T$~vs.~$\eta$
correlation. This validates: ROOT integration, library packaging,
\texttt{SetupCAP.sh} environment, the Pythia~8 link via the davidchall/hep
tap, and the CMake install layout.

\noindent
The runner GUI exposes the macro as a one-click \emph{Smoke test (proof of
life)} button.

% =============================================================================
\section{Documentation added}
% =============================================================================

\begin{itemize}[leftmargin=2em,itemsep=0pt]
  \item \texttt{README.md} --- short quick-start.
  \item \texttt{BUILDING.md} --- build modes, run modes, troubleshooting,
        Wayne~State Grid notes.
  \item \texttt{CODEBASE\_MAP.md} --- complete architectural reference for
        every module in \texttt{src/} (and an inventory of the vendored
        FastJet library).
  \item \texttt{JETS\_KNOWN\_ISSUES.md} --- the Jets-module bug list with
        file/line references.
  \item \texttt{MISSING\_CLASSES.md} --- the missing-class punch-list.
  \item \texttt{TECHNICAL\_REPORT.tex} --- this document.
\end{itemize}

% =============================================================================
\section{Recommended next steps}
% =============================================================================

\begin{enumerate}[leftmargin=2em,itemsep=0.3em]
  \item \textbf{Implement the trivial missing classes first.} \texttt{RunAnalysis}
        is little more than \texttt{class RunAnalysis : public Task \{\};}
        plus the \texttt{ClassDef}/\texttt{ClassImp} macros; with it, the
        \texttt{.ini} load path completes and the next missing class
        becomes visible. \texttt{ParticleTypeTask} (DB loader) is the
        natural second target.

  \item \textbf{Decide on the configuration-key migration policy.} Options:
        (a) extend the back-compat shim from Section~\ref{sec:taskpatch}
        to per-task keys (lots of mechanical edits, low risk),
        (b) rewrite the project \texttt{.ini} files to the new
        \texttt{<typeName>}-aware structure (cleaner but requires
        per-task knowledge of which type-name each task uses),
        (c) accept both syntaxes via a centralised helper in
        \texttt{Configuration} or \texttt{ConfigurationManager} that
        applies a translation table at lookup time.

  \item \textbf{Repair or delete the Jets module.} The list in
        \texttt{JETS\_KNOWN\_ISSUES.md} is comprehensive. Until a maintainer
        chooses to engage it, the build correctly keeps it disabled.

  \item \textbf{Repair \texttt{PythiaEventReader.cpp} or remove it.} It is
        currently excluded from the build. If the TTree-reading workflow it
        embodied is still wanted, the schema needs to be re-derived; if
        not, the file should be deleted.

  \item \textbf{Adopt the new build/install/runner tooling.} The
        \texttt{./install} entry-point is now the recommended user-facing
        command. CI scripts can use \texttt{./install --headless}; users
        with a desktop environment use \texttt{./install --gui}; HPC head
        nodes use \texttt{./install --cli}. None of these depend on the
        missing classes.
\end{enumerate}

% =============================================================================
\appendix
\section{Files added or modified}
% =============================================================================

\subsection*{Added}

\begin{itemize}[leftmargin=2em,itemsep=0pt]
  \item \texttt{CMakeLists.txt} (top-level)
  \item \texttt{cmake/FindPythia8.cmake}
  \item \texttt{cmake/FindFastJet.cmake}
  \item \texttt{install} (launcher)
  \item \texttt{setup-cap} (CLI installer)
  \item \texttt{setup-cap-gui} (Tk installer)
  \item \texttt{run-cap} (Tk job runner)
  \item \texttt{scripts/cap-logging.sh}
  \item \texttt{src/Macros/SmokeTest.C}
  \item \texttt{README.md}, \texttt{BUILDING.md},
        \texttt{CODEBASE\_MAP.md}, \texttt{JETS\_KNOWN\_ISSUES.md},
        \texttt{MISSING\_CLASSES.md}, \texttt{TECHNICAL\_REPORT.tex}
\end{itemize}

\subsection*{Modified}

\begin{itemize}[leftmargin=2em,itemsep=0pt]
  \item \texttt{SetupCAP.sh} (full rewrite for portability)
  \item \texttt{src/CMakeLists.txt} (modernised)
  \item \texttt{src/Base/Task.cpp} (back-compat shim, Section~\ref{sec:taskpatch})
  \item \texttt{src/CAPPythia/CMakeLists.txt} (excludes the broken reader)
  \item \texttt{src/CAPPythia/CAPPythiaLinkDef.h} (excludes the broken reader)
  \item \texttt{src/CAPPythia/PythiaEventGenerator.cpp} (twelve repair sites)
  \item \texttt{src/Glauber/GlauberHistos.cpp} (return statement, unused vars)
  \item \texttt{src/Glauber/NuclearDistribution.cpp} (unused var)
  \item \texttt{src/Helpers/CanvasList.cpp} (unused var)
  \item \texttt{src/Jets/CMakeLists.txt} (uses \texttt{Pythia8::Pythia8} target)
  \item \texttt{src/Jets/*.{hpp,cpp}}
        (\texttt{ParticleTypeList} $\to$ \texttt{ParticleDb}, typo fix)
  \item All twenty-two \texttt{src/<Module>/CMakeLists.txt}
        (\texttt{\$ENV\{CAP\_LIB\_PATH\}} $\to$ \texttt{\$\{CAP\_INSTALL\_LIBDIR\}})
  \item \texttt{.gitignore} (logs, .cap-config, in-tree pythia/fastjet builds)
\end{itemize}

\end{document}
